% Generated draft: requires mathematical and source review.
\documentclass[11pt]{article}
\usepackage[margin=25mm]{geometry}
\usepackage{fontspec}
\setmainfont{Latin Modern Roman}
\usepackage{amsmath,amssymb,mathtools,unicode-math}
\setmathfont{Latin Modern Math}
\usepackage{enumitem,needspace,xurl,hyperref}
\setlength{\parindent}{0pt}
\setlength{\parskip}{5pt}
\setlength{\emergencystretch}{3em}
\newcommand{\N}{\mathbb{N}}
\newcommand{\Z}{\mathbb{Z}}
\newcommand{\Q}{\mathbb{Q}}
\newcommand{\R}{\mathbb{R}}
\newcommand{\C}{\mathbb{C}}
\newcommand{\F}{\mathbb{F}}
\newcommand{\A}{\mathbb{A}}
\newcommand{\PP}{\mathbb{P}}
\newcommand{\E}{\mathbb{E}}
\newcommand{\eps}{\varepsilon}
\newcommand{\dd}{\,\mathrm{d}}
\begin{document}
\section*{UnsolvedMath ID: 20001909}
\textbf{Source problem code:} catalogue entry
\subsection{Definitions and mathematical statement}
Let $(M,g)$ be a smooth closed connected oriented surface of genus $q\ge2$, with area form $\dd A_g$, Gaussian curvature $K_g$, and magnetic form $\Omega=f\,\dd A_g$. Fix $k>0$. With $J$ rotation by $\pi/2$, choose the Lorentz-force convention
\[
 \nabla_t\dot\gamma=f(\gamma)J\dot\gamma,
 \qquad \tfrac12|\dot\gamma|_g^2=k.
\]
Rescaling time to unit speed identifies the energy surface with $SM$ and gives the vector field $F_k=X+\lambda_kV$, where $X$ is the geodesic vector field, $V$ generates positive fiber rotations and $\lambda_k=f/\sqrt{2k}$. Its transverse magnetic Jacobi equation is
\[
 y''+Q_k(\gamma,u)y=0,
 \qquad Q_k(x,u)=K_g(x)-\frac{\dd f_x(Ju)}{\sqrt{2k}}
                   +\frac{f(x)^2}{2k}.
\]
Here absence of conjugate points means that a nonzero solution of this reduced transverse equation cannot vanish at two distinct times along any orbit. The question concerns the whole positive energy surface, not one orbit.

\textbf{Question 7.4.1, attributed to Paternain.} If this magnetic flow has zero topological entropy and has no conjugate points, must $K_g$ be constant and $f$ be constant? This is precisely the source's meaning of local symmetry: $g$ has constant curvature and $\Omega$ is a constant multiple of its area form. Positive constant rescaling of time does not affect the condition of zero entropy.

\textbf{Supplied partial result, requiring verification.} The research aid asserts that the additional hypothesis $Q_k\le0$ everywhere yields $Q_k\equiv0$ and hence
\[
 f\equiv c,\qquad K_g\equiv-\frac{c^2}{2k}.
\]
Since $\chi(M)<0$, this model has $c\ne0$ and the normalized flow is of horocycle type. It further asserts the following Hopf defect statement under no conjugate points alone. Use the invariant measure $\dd\mu=\dd A_g\,\dd\theta/(2\pi)$ on $SM$; its total mass is $\operatorname{Area}_g(M)$, not one. For a measurable Green--Riccati solution $r$ of $F_kr+r^2+Q_k=0$, the claimed identity is
\[
 \int_{SM}r^2\,\dd\mu
 =D_k:=-2\pi\chi(M)-\frac1{2k}\int_M f^2\,\dd A_g\ge0.
\]
Thus $k\ge \bigl(\int_M f^2\,\dd A_g\bigr)/(-4\pi\chi(M))$, and the aid identifies $D_k=0$ with the same constant-curvature, constant-field model. These are special-case and equality claims; neither supplies the missing implication from zero entropy to zero defect under only the original hypotheses. The Riccati existence/integrability and entropy arguments need mathematical review.

\textbf{Adjacent source material.} The extract says that earlier questions can also be asked at a fixed, possibly sufficiently high, magnetic energy. It reports that [Bialy2000] proves an analogue of Theorem 7.0.2 under a conformal-flatness assumption, conjecturally unnecessary. Neither that theorem nor its ambient hypotheses is supplied here; no precise new version is inferred.

The extract also contains the separate Section 8.1 question: does every closed manifold of dimension at least two admit a Riemannian or Finsler metric whose geodesic flow is ergodic? Under the usual interpretation, ergodicity is with respect to the normalized natural invariant Liouville measure on the unit tangent/cotangent bundle. The extract reports positive cases for surfaces, three-manifolds, products of factors of dimension at most three, and spheres; it also reports that every embedded surface admits arbitrarily small $C^0$ perturbations to embedded surfaces with ergodic induced geodesic flow. It attributes these constructions to Donnay's focusing caps and warns that the general question remains open in that source despite contrary reports in Berger. These are historical source claims, not a current literature assessment. The final heading ``8.2. Positive curvature'' contains no further question.
\subsection{Short English statement}
Paternain asks whether a zero-entropy magnetic flow without conjugate points on a closed surface of genus at least two must have constant curvature and constant magnetic field. The supplied aid adds a magnetic-curvature sign hypothesis or a zero Hopf-defect hypothesis; it does not settle the stated question. The merged extract also asks for ergodic geodesic metrics on arbitrary closed manifolds.
\subsection{Sources}
\begin{itemize}
\item Local UnsolvedMath registry, permanent ID 20001909, source code catalogue entry; canonical location \texttt{aim-geometry-notes.json}, notes[246].
\item American Institute of Mathematics, Geodesics problem list, Question 7.4.1 (Paternain), followed by Section 8.1: \url{https://aimath.org/pastworkshops/geodesicsproblems.pdf}. The record explicitly warns of text captured from the next section.
\item Supplied magnetic-flow reference label [Bialy2000], with the missing cross-reference Theorem 7.0.2. No full bibliographic entry is supplied.
\item Supplied Section 8.1 labels: [Donnay1988II], [Katok1994], [Bur-Ger1994], [Bur-Ged], [Don-Pugh2004], and [Donnay1988I]; contrary-solution reports are located at p.~87 of [Berger2000] and Section 10.9 of [Berger2003]. Full entries are absent.
\item Registry research attempt 1, dated 2026-08-14, labelled partial result, rigor ``proved'', novelty confidence low. It is explicitly a machine-generated, non-peer-reviewed aid to be independently verified before citation.
\end{itemize}
\end{document}
